\part{System Architecture}

\chapter{System Architecture}

\section{Architectural Drivers}

The architecture was not chosen in a vacuum. The most important design drivers were ranked before the
service boundaries were fixed because different priorities would have produced a meaningfully different
system.

\begin{table}[H]
\centering
\caption{Ranked architectural drivers}
\label{tab:architectural-drivers}
\begin{tabularx}{\textwidth}{c p{4.2cm} X}
\toprule
\textbf{Rank} & \textbf{Driver} & \textbf{Why it mattered} \\
\midrule
1 & Collaboration correctness and convergence & The product fails if two users cannot trust what they are seeing while editing the same document. This pushed the design toward a CRDT-backed push model rather than naive polling or last-write-wins merging. \\
2 & Low-latency user experience & Typing, presence, and AI interaction feedback must feel responsive enough to resemble a real collaborative editor, not a batch workflow. \\
3 & Maintainable modular codebase & The assignment explicitly evaluates repository organization and feature decomposition. The architecture had to support multiple owners working in parallel without turning into a god-service monolith. \\
4 & AI safety, predictability, and cost control & AI is a core product feature, not a novelty. This drove the asynchronous job model, scoped prompting, and explicit policy/quota controls. \\
5 & Demo and reproducibility reliability & The MVP had to be runnable under grading pressure. This drove root orchestration scripts, SQLite portability, Docker Compose, CI, and deterministic reset behavior. \\
\bottomrule
\end{tabularx}
\end{table}

These rankings directly shaped the implementation: the realtime service is separated from the REST API,
AI uses asynchronous job handling rather than inline blocking calls, persistence is centralized in
SQLite for deterministic setup, and the repository is organized by service and domain boundaries rather
than by one large application file.

\section{C4 Level 1 --- System Context}

Figure~\ref{fig:c4-l1} shows the system at the outermost level. The collaborative editor interacts with
human collaborators, a target identity provider, and an LLM provider. This level answers the question:
what is inside the system boundary, and what external actors or services affect behavior?

\begin{figure}[H]
  \centering
  \includegraphics[width=0.82\textwidth]{c4-level-1-system-context.png}
  \caption{C4 Level 1 --- System context diagram}
  \label{fig:c4-l1}
\end{figure}

\implementationnote{The delivered MVP uses a visible local login flow backed by
\texttt{POST /auth/login} rather than a full external OAuth/OIDC identity provider. The system-context
diagram therefore captures the target production-style identity boundary, while the MVP demonstrates the
same trust boundary with a simpler bearer-token implementation suitable for coursework and demos.}

\section{C4 Level 2 --- Container Diagram}

Figure~\ref{fig:c4-l2} decomposes the system into the major runtime containers. The frontend is the user
entry point. The backend API owns document lifecycle, permissions, versioning, export, AI orchestration,
and session issuance. The realtime service is dedicated to Yjs synchronization, awareness, and role-aware
write gating. The AI integration service abstracts provider selection, and the database stores the shared
system state.

\begin{figure}[H]
  \centering
  \includegraphics[width=0.9\textwidth]{c4-level-2-container.png}
  \caption{C4 Level 2 --- Container diagram}
  \label{fig:c4-l2}
\end{figure}

\implementationnote{The delivered MVP uses SQLite rather than a production-scale database deployment.
The architecture keeps the same container responsibilities, but the concrete persistence technology is
optimized for a reproducible academic environment. The realtime service and backend share the same SQLite
state to keep sessions, permissions, versions, and backend-driven collaboration events aligned.}

\section{C4 Level 3 --- Backend API Component Diagram}

The backend was chosen as the container to zoom into because it contains most of the assignment's
cross-cutting concerns: document CRUD, RBAC, AI job orchestration, export job orchestration, versioning,
and realtime session issuance. Figure~\ref{fig:c4-l3} shows the controller/service/repository split used
to keep the API layer maintainable.

\begin{figure}[H]
  \centering
  \includegraphics[width=\textwidth]{c4-level-3-backend-components.png}
  \caption{C4 Level 3 --- Backend API component diagram}
  \label{fig:c4-l3}
\end{figure}

\implementationnote{The delivered backend closely follows this structure: routes/controllers handle HTTP
translation, auth middleware validates bearer tokens, services encode document/AI/export/session rules,
repositories isolate SQL, and the realtime service is notified through document-scoped events rather than
through direct in-process mutation. The MVP uses SQLite and local bearer-token authentication in place of
the fuller target stack shown in the conceptual diagram.}

\section{Feature Decomposition}

The architecture is organized around features/modules rather than a single monolithic application file.
Each module exposes a narrow interface, depends on a small set of collaborators, and can be tested in
relative isolation.

\begin{table}[H]
\centering
\caption{Feature decomposition}
\label{tab:feature-decomposition}
\begin{tabularx}{\textwidth}{p{3.2cm}p{4.3cm}X}
\toprule
\textbf{Module} & \textbf{Primary Responsibility} & \textbf{Dependencies / Exposed Interface} \\
\midrule
Frontend editor and state & Render document content, local editing state, AI review panel, permissions, history, export controls & Depends on Backend API and Realtime Service; exposes user actions as REST and websocket calls \\
Realtime synchronization layer & Relay Yjs sync and awareness frames, enforce role-aware writes, react to backend permission/revert events & Depends on signed sessions, shared DB state, and frontend websocket clients; exposes session-ready and live sync behavior \\
Backend API layer & Authenticate users, enforce RBAC, manage document lifecycle, issue realtime sessions, orchestrate AI and exports & Depends on repositories, AI provider adapter, and realtime-facing event model; exposes REST contracts and standard error schema \\
AI integration service & Translate user AI requests into provider calls, manage async jobs, unify error handling & Depends on provider implementation and backend persistence; exposes job creation and polling semantics \\
Document storage and versioning & Persist documents, versions, policies, jobs, permissions, idempotency keys, and audit records & Depends on SQLite; exposes repository methods to backend services \\
Shared contract/utilities layer & Session-token helpers and reusable cross-service logic & Depends on no single service; exposes shared helpers used by backend and realtime \\
\bottomrule
\end{tabularx}
\end{table}

\section{AI Integration Design}

The AI assistant is treated as a core product feature, not a cosmetic API add-on. That means the design
must answer four questions simultaneously: what context the model sees, how users review and apply
results, how collaboration behaves while the AI is running, and how cost and failure are controlled.

\subsection{Context and Scope}

The delivered MVP sends the selected text plus a small amount of surrounding context to the AI provider.
This is the most pragmatic choice for coursework: it is cheaper than sending the full document, keeps
latency lower, and aligns with how rewrite/summarize/translate are actually invoked in the UI. The
backend contract requires the frontend to send a client-side selection snapshot containing the selected
range, the selected text, and the nearby context so that AI jobs can be linked to a concrete document
state.

For very long documents, the system intentionally does not default to sending the entire document. This
keeps token usage bounded and reduces privacy exposure. A future production design could add
chunked/restructure flows, but the delivered MVP focuses on scoped selection-based tasks that are easier
to reason about and test.

\subsection{Suggestion UX}

AI suggestions are presented in a side panel rather than silently modifying the document. This preserves
user control and matches the assignment's emphasis on AI as a product feature rather than a hidden
automation step. The delivered MVP supports:

\begin{itemize}
  \item full apply of the AI output,
  \item partial apply by selecting only a portion of the AI output,
  \item explicit rejection without modifying document content,
  \item immediate visibility of AI job state.
\end{itemize}

Applying AI is treated as a document-changing action and therefore creates version snapshots for
traceability. This design also aligns with the version history and revert model: AI is not special-cased
 into an untracked side path.

\subsection{AI During Collaboration}

The design intentionally avoids hard-locking a region when AI is invoked. Instead, the system snapshots
the selected text, lets the AI run asynchronously, and allows other collaborators to continue editing.
When the AI response arrives, the frontend presents it as a proposal. If the relevant text has changed in
the meantime, the user must explicitly decide whether to apply or discard the AI output. This is a
reconcile-first design rather than a lock-first design.

This choice is justified by the architectural drivers. Locking would be simpler to explain, but it would
damage collaboration fluidity and create awkward waiting states. Reconciliation with explicit user review
keeps the collaboration model more honest, even if it requires careful versioning and conflict-safe save
behavior.

\subsection{Prompt Design and Provider Strategy}

Prompts are task-specific and template-based. The frontend chooses the user-facing action
(rewrite/summarize/translate), while the backend/provider layer is responsible for constructing the final
provider request. This keeps prompt construction out of the browser and allows prompt logic or provider
selection to evolve without changing the frontend contract.

The delivered MVP uses a provider abstraction with two concrete modes:

\begin{itemize}
  \item \texttt{stub}: a deterministic provider used for reliable demos and tests,
  \item \texttt{lmstudio}: a local OpenAI-compatible provider used when a real local model endpoint is
  available.
\end{itemize}

This is a pragmatic compromise between the architectural goal of AI usefulness and the assignment's need
for repeatable demonstrations.

\subsection{Async Job Model}

AI operations are modeled as asynchronous jobs. The backend validates access, policy, quota, and input
shape, creates a job record, and returns quickly with \texttt{jobId}, \texttt{statusUrl}, and an initial
status. The frontend polls until the job reaches \texttt{SUCCEEDED} or \texttt{FAILED}. This model avoids
blocking request threads and gives the client a precise way to distinguish “still running” from “failed”
and “succeeded.”

\section{API Design}

The API style follows a hybrid model:

\begin{itemize}
  \item \textbf{REST} for CRUD, permissions, versions, exports, AI jobs, and session issuance,
  \item \textbf{WebSocket + Yjs} for low-latency collaborative state propagation,
  \item \textbf{Polling} from the frontend only for long-running AI jobs and export jobs.
\end{itemize}

This division is deliberate. REST keeps persistence-oriented operations predictable and easy to test.
Realtime collaboration is handled by a dedicated push channel because document editing is inherently a
shared live interaction. Long-running AI and export operations fit naturally into job-based REST flows.

\subsection{Core API Contracts}

The assignment brief requires at least document CRUD, realtime session management, AI invocation, and
user/permission management. The delivered backend provides these contracts concretely and documents them
through OpenAPI and contract notes in the repository.

\begin{table}[H]
\centering
\caption{Core REST contract families used by the MVP}
\label{tab:core-api-families}
\begin{tabularx}{\textwidth}{p{3.8cm}X}
\toprule
\textbf{API family} & \textbf{Concrete delivered endpoints} \\
\midrule
Authentication & \texttt{POST /auth/login}, \texttt{GET /health} \\
Document lifecycle & \texttt{POST /documents}, \texttt{GET /documents/:documentId}, \texttt{PUT /documents/:documentId/content} \\
Versioning & \texttt{GET /documents/:documentId/versions}, \texttt{POST /documents/:documentId/revert} \\
Permissions and AI policy & \texttt{GET/PUT/DELETE /documents/:documentId/permissions}, \texttt{GET/PUT /documents/:documentId/ai-policy} \\
Realtime session management & \texttt{POST /sessions} \\
AI jobs & \texttt{POST /ai/rewrite}, \texttt{POST /ai/summarize}, \texttt{POST /ai/translate}, \texttt{GET /ai/jobs/:jobId} \\
Export jobs & \texttt{POST /documents/:documentId/export}, \texttt{GET /exports/:jobId}, \texttt{GET /exports/:jobId/download} \\
\bottomrule
\end{tabularx}
\end{table}

All non-2xx responses use a shared JSON error shape:

\begin{quote}
\texttt{\{ "error": \{ "code": "<CODE>", "message": "<MESSAGE>", "details": \{ ...optional... \} \} \}}
\end{quote}

This is important because the assignment brief explicitly asks how the client can distinguish between AI
latency, AI failure, quota issues, permission problems, and not-found conditions. The standard error
schema gives the frontend a uniform contract for all of those cases.

Representative delivered contracts are shown below.

\paragraph{Document create}
\begin{lstlisting}[language={}]
POST /documents

Request JSON:
{ "title": "Project brief", "content": "Initial content" }

Response JSON (201):
{
  "documentId": "doc_123",
  "title": "Project brief",
  "ownerId": "user_1",
  "createdAt": "ISO-8601",
  "updatedAt": "ISO-8601",
  "currentVersionId": "ver_1"
}
\end{lstlisting}

\paragraph{Document load}
\begin{lstlisting}[language={}]
GET /documents/:documentId

Response JSON (200):
{
  "documentId": "doc_123",
  "title": "Project brief",
  "content": "Latest text",
  "updatedAt": "ISO-8601",
  "currentVersionId": "ver_3",
  "role": "editor",
  "revisionId": "rev_3"
}
\end{lstlisting}

\paragraph{Realtime session issuance}
\begin{lstlisting}[language={}]
POST /sessions

Request JSON:
{ "documentId": "doc_123" }

Response JSON (200):
{
  "sessionId": "sess_123",
  "wsUrl": "ws://localhost:3001/ws?token=...",
  "role": "viewer|editor|owner"
}
\end{lstlisting}

\paragraph{AI rewrite invocation}
\begin{lstlisting}[language={}]
POST /ai/rewrite

Request JSON:
{
  "documentId": "doc_123",
  "selection": { "start": 0, "end": 25 },
  "selectedText": "Original selected text",
  "contextBefore": "",
  "contextAfter": " remaining content",
  "instruction": "Make this more concise",
  "baseVersionId": "ver_3",
  "requestId": "req_ai_001"
}

Response JSON (202):
{
  "jobId": "aijob_123",
  "statusUrl": "/ai/jobs/aijob_123",
  "status": "PENDING"
}
\end{lstlisting}

\paragraph{AI job polling}
\begin{lstlisting}[language={}]
GET /ai/jobs/:jobId

Response JSON (200):
{
  "jobId": "aijob_123",
  "statusUrl": "/ai/jobs/aijob_123",
  "status": "SUCCEEDED",
  "proposedText": "Rewritten text",
  "baseVersionId": "ver_3",
  "createdAt": "ISO-8601",
  "updatedAt": "ISO-8601"
}
\end{lstlisting}

\paragraph{Permission grant}
\begin{lstlisting}[language={}]
PUT /documents/:documentId/permissions

Request JSON:
{ "requestId": "req_perm_001", "targetUserId": "user_2", "role": "editor" }

Response JSON (200):
{ "documentId": "doc_123", "targetUserId": "user_2", "role": "editor", "updatedAt": "ISO-8601" }
\end{lstlisting}

\subsection{Authentication and Authorization}

The target architecture includes an external identity provider, but the delivered MVP uses an explicit
local login flow backed by bearer tokens. The backend nevertheless enforces the same high-level concerns
that would remain in a production system: authentication gates protected resources, authorization is
document-scoped, and the UI is not treated as the source of truth.

\begin{table}[H]
\centering
\caption{Document-scoped role matrix in the delivered MVP}
\label{tab:role-matrix}
\begin{tabularx}{\textwidth}{>{\raggedright\arraybackslash}X c c c}
\toprule
\textbf{Action} & \textbf{Viewer} & \textbf{Editor} & \textbf{Owner} \\
\midrule
View/open document & Yes & Yes & Yes \\
Receive realtime updates/presence & Yes & Yes & Yes \\
Send collaborative edit operations & No & Yes & Yes \\
Persist content via REST save/autosave & No & Yes & Yes \\
View version history & Yes & Yes & Yes \\
Revert to previous version & No & No & Yes \\
Export document & Yes & Yes & Yes \\
Grant/update collaborator roles & No & No & Yes \\
Revoke access & No & No & Yes \\
Invoke AI rewrite/summarize/translate & No & Yes & Yes \\
Configure AI policy and quota & No & No & Yes \\
\bottomrule
\end{tabularx}
\end{table}


\implementationnote{The delivered system fully demonstrates owner/editor/viewer behavior. “Admin” is
kept as a design-level role in the report because it remains meaningful for future policy and
organization-level control, but the MVP's primary demonstrated management path is document ownership.}

\subsection{Communication Model}

The communication model is intentionally hybrid.

\begin{itemize}
  \item \textbf{Document open flow.} The frontend first loads a document snapshot through REST and then
  requests a signed realtime session URL from the backend.
  \item \textbf{Live collaboration.} The realtime service validates the signed session token and relays Yjs
  sync and awareness updates among collaborators.
  \item \textbf{Disconnect behavior.} If realtime becomes unavailable, the frontend falls back to offline/manual-save mode.
  \item \textbf{AI completion.} AI uses async jobs with polling instead of websocket completion events in
  the MVP because the job semantics are easier to reason about and test.
\end{itemize}

Additional sequence-diagram artifacts for document-open/realtime-join and AI-job execution are included
in the repository's editable Mermaid sources, but omitted from the main PDF to keep the report focused on
the assignment's required C4 and ER views.

\begin{figure}[H]
  \centering
  \includegraphics[width=0.95\textwidth]{c4-level-2-container.png}
  \caption{Runtime containers involved in document open and collaboration sync}
  \label{fig:seq-open-sync}
\end{figure}

\implementationnote{The delivered system uses Yjs over WebSocket for text synchronization and presence.
Polling is used only for async AI jobs and export jobs, and internally for backend-driven permission/
revert event delivery to realtime rooms. It is not used to poll text diffs every second.}

\section{Code Structure and Repository Organization}

The final repository intentionally uses a monorepo. For a four-person academic team with strongly
connected frontend, backend, realtime, and AI layers, a monorepo reduces coordination overhead and makes
contract drift easier to detect.

\subsection{Top-Level Layout}

\begin{lstlisting}[basicstyle=\ttfamily\small]
collaborative-editor-ai/
+-- frontend/        React + TypeScript client
+-- backend/         Express API, SQLite persistence, RBAC, AI jobs, export
+-- realtime/        WebSocket + Yjs relay service
+-- ai-service/      AI provider abstraction and prompt execution
+-- shared/          shared contracts such as session token signing
+-- docs/            architecture notes, Mermaid sources, diagram exports
+-- scripts/         root install/dev/test/demo orchestration
+-- tests/           root integration smoke tests
`-- final-report/    final written submission package
\end{lstlisting}

\subsection{Shared Code and Config}

The shared layer is intentionally narrow. It contains only cross-service helpers that justify shared
ownership, such as session-token signing/verification. This avoids over-abstracting everything into a
shared package while still keeping critical contracts consistent. Configuration is managed through
environment variables and \texttt{.env.example} files; secrets are excluded from version control.

\subsection{Testing Structure}

The repository uses layered testing rather than relying on only one type of test.

\begin{itemize}
  \item Backend: unit/integration tests with Node's built-in test runner
  \item Frontend: component/integration-style tests with Vitest
  \item Realtime: websocket/Yjs integration tests with Node's built-in test runner
  \item Root: a smoke integration test that starts backend and realtime together and validates wiring
\end{itemize}

This structure mirrors the feature decomposition: REST, realtime, and UI are each tested in their own
layer, and the root smoke test catches cross-service drift.

\section{Data Model}

Figure~\ref{fig:er-model} shows the core entity relationships. The design centers on documents, users,
permissions, versions, and AI interactions, but the delivered MVP also persists AI policy, export jobs,
realtime events, audit logs, and idempotency keys to make the runtime predictable and observable.

\begin{figure}[H]
  \centering
  \includegraphics[width=0.96\textwidth]{er-data-model.png}
  \caption{Entity-relationship diagram for the collaborative editor domain}
  \label{fig:er-model}
\end{figure}

\implementationnote{The delivered schema uses SQLite and includes implementation-supporting entities
such as AI policy, export jobs, realtime events, audit logs, and idempotency keys in addition to the
core design entities. This keeps the delivered MVP honest about the data required to support jobs,
permissions, versioning, and collaboration behavior.}

\section{Architecture Decision Records}

\subsection*{ADR-001 --- Realtime consistency strategy}
\textbf{Status:} Accepted.\\
\textbf{Context:} The system needs simultaneous editing, reconnect safety, and convergence without the
team implementing a custom operational transform algorithm.\\
\textbf{Decision:} Use a CRDT-based collaboration approach built on Yjs, with the realtime service acting
as a relay and permission gate rather than a custom transform engine.\\
\textbf{Consequences:} This improves convergence and offline/reconnect resilience, but requires careful
handling of room lifecycle, awareness state, and persistence boundaries.\\
\textbf{Alternatives considered:} Custom OT server logic, naive last-write-wins, and polling-based text
sync were rejected due to complexity or poor correctness.

\subsection*{ADR-002 --- API style and async job model}
\textbf{Status:} Accepted.\\
\textbf{Context:} The system mixes CRUD-like document operations with long-running AI and export actions.\\
\textbf{Decision:} Use REST for document, permission, version, AI job, export job, and session endpoints;
use WebSocket for realtime collaboration; and model AI/export as asynchronous jobs.\\
\textbf{Consequences:} The API is concrete and testable, but the frontend must poll job status and the
backend must persist job state.\\
\textbf{Alternatives considered:} synchronous AI calls, GraphQL for all interactions, and polling for live
collaboration were rejected.

\subsection*{ADR-003 --- Repository structure strategy}
\textbf{Status:} Accepted.\\
\textbf{Context:} The team needed a code structure that matched a multi-service system while minimizing
coordination overhead for a small academic group.\\
\textbf{Decision:} Use a monorepo with service directories and keep persistence logic grouped by domain
repositories rather than a single giant application repository file.\\
\textbf{Consequences:} Shared contracts are easier to keep aligned, but discipline is required to prevent
tight coupling.\\
\textbf{Alternatives considered:} separate repositories and hybrid split-repo models were rejected due to
contract drift and coordination cost.

\subsection*{ADR-004 --- AI provider strategy}
\textbf{Status:} Accepted.\\
\textbf{Context:} The system needs rewrite, summarize, and translate behaviors that are useful during a
demo but not dependent on an always-available external cloud provider.\\
\textbf{Decision:} Keep AI behind a provider abstraction supporting a deterministic stub mode and a local
LM Studio/OpenAI-compatible mode.\\
\textbf{Consequences:} Demo reliability improves and tests remain deterministic, but the system is still
architected so a different provider can be added later.\\
\textbf{Alternatives considered:} hard-wiring one hosted provider or requiring self-hosted production-scale
model infrastructure for the MVP.
